\documentclass[11pt,a4paper]{article}

% ── Codificación, fuentes y lenguaje ───────────────────────────────────────────
\usepackage[utf8]{inputenc}
\usepackage[T1]{fontenc}
\usepackage[default,scale=0.95]{sourcesanspro}
\renewcommand{\familydefault}{\sfdefault}
\usepackage[spanish,es-noshorthands]{babel}

% ── Geometría y espaciado ─────────────────────────────────────────────────────
\usepackage[top=2.2cm, bottom=2.2cm, left=2.2cm, right=2.2cm]{geometry}
\usepackage{setspace}
\setstretch{1.15}
\usepackage{parskip}

% ── Matemáticas y unidades ────────────────────────────────────────────────────
\usepackage{amsmath}
\usepackage{amssymb}
\usepackage{siunitx}

% ── Circuitos ─────────────────────────────────────────────────────────────────
\usepackage[european, straightvoltages]{circuitikz}

% ── Tablas ────────────────────────────────────────────────────────────────────
\usepackage{booktabs}
\usepackage{tabularx}
\usepackage{array}
\usepackage{longtable}
\usepackage{colortbl}
\usepackage{enumitem}

% ── Colores, cajas e iconos ───────────────────────────────────────────────────
\usepackage[dvipsnames]{xcolor}
\usepackage{tcolorbox}
\tcbuselibrary{skins, breakable}
\usepackage{fontawesome5}

% ── Secciones con barra de color (infografía) ───────────────────────────────
\usepackage{titlesec}
\titleformat{\section}
  {\normalfont\LARGE\bfseries\color{azulTitulo}}
  {\thesection}{1em}{}
  [\color{bandaAzul}\rule{\linewidth}{1.6pt}]
\titlespacing*{\section}{0pt}{22pt}{10pt}
\titleformat{\subsection}
  {\normalfont\large\bfseries\color{azulTitulo}}
  {\thesubsection}{1em}{}
\titlespacing*{\subsection}{0pt}{14pt}{6pt}

% ── Cabeceras y pies de página ────────────────────────────────────────────────
\usepackage{fancyhdr}
\usepackage{lastpage}
\pagestyle{fancy}
\fancyhf{}

% ── Hiperenlaces ──────────────────────────────────────────────────────────────
\usepackage[colorlinks=true, linkcolor=azulTitulo, urlcolor=azulTitulo]{hyperref}
\sloppy
\emergencystretch 3em

% ── Paleta de colores de la marca ─────────────────────────────────────────────
\definecolor{azulTitulo}{HTML}{1B2A6B}
\definecolor{bandaAzul}{HTML}{1F3A93}
\definecolor{azulClaro}{HTML}{2E5A9C}
\definecolor{cyanAcento}{HTML}{0E7C86}
\definecolor{naranjaAcento}{HTML}{E07B39}
\definecolor{verdeAcento}{HTML}{1A6B2F}
\definecolor{rojoAcento}{HTML}{C0392B}
\definecolor{grisFondo}{HTML}{F5F6FA}
\definecolor{grisBorde}{HTML}{C9CFE0}

% ── Tarjetas de datos (infografía) ────────────────────────────────────────────
\tcbset{
  tarjetaDato/.style={
    enhanced, breakable,
    colback=azulClaro!6, colframe=azulClaro,
    fonttitle=\bfseries, coltitle=white,
    attach boxed title to top left={yshift=-2mm, xshift=4mm},
    boxed title style={colback=azulClaro, sharp corners},
    top=6pt, bottom=6pt, left=8pt, right=8pt,
  },
  tarjetaIcono/.style={
    enhanced, breakable, sharp corners,
    colback=white, colframe=azulClaro, boxrule=0.8pt,
    fonttitle=\bfseries, coltitle=azulTitulo,
    top=4pt, bottom=4pt, left=8pt, right=8pt,
  },
  cajaTitulo/.style={enhanced, breakable,
    colback=azulTitulo!8, colframe=azulTitulo,
    fonttitle=\bfseries\large, coltitle=white,
    attach boxed title to top left={yshift=-2mm, xshift=4mm},
    boxed title style={colback=azulTitulo},
    top=6pt, bottom=6pt, left=8pt, right=8pt},
  cajaContenido/.style={enhanced, breakable,
    colback=grisFondo, colframe=grisBorde,
    fonttitle=\bfseries, coltitle=azulTitulo,
    top=4pt, bottom=4pt, left=8pt, right=8pt},
  cajaFortaleza/.style={enhanced, breakable,
    colback=verdeAcento!8, colframe=verdeAcento,
    fonttitle=\bfseries, coltitle=verdeAcento,
    top=4pt, bottom=4pt, left=8pt, right=8pt},
  cajaMejora/.style={enhanced, breakable,
    colback=naranjaAcento!8, colframe=naranjaAcento,
    fonttitle=\bfseries, coltitle=naranjaAcento,
    top=4pt, bottom=4pt, left=8pt, right=8pt},
  cajaRecomendacion/.style={enhanced, breakable,
    colback=grisFondo, colframe=grisBorde,
    fonttitle=\bfseries, coltitle=azulTitulo,
    top=4pt, bottom=4pt, left=8pt, right=8pt},
}

% ── Banda de título: apertura de documento tipo infografía ─────────────────────
\newcommand{\bandaTitulo}[2]{%
  \begin{tcolorbox}[enhanced, colback=bandaAzul, colframe=bandaAzul,
    boxrule=0pt, arc=0pt, outer arc=0pt, width=\linewidth,
    halign=center, valign=center,
    top=14pt, bottom=14pt, left=10pt, right=10pt]
    {\LARGE\bfseries\color{white}#1}\par\vspace{4pt}%
    {\small\color{white!78!black}#2}%
  \end{tcolorbox}\vspace{8pt}}

% ── Ícono + texto (encabezados de sección y elementos) ────────────────────────
\newcommand{\iconotexto}[2]{\textcolor{azulTitulo}{\faIcon{#1}}~#2}

% ─────────────────────────────────────────────────────────────────────────────

\fancyhead[L]{\small\color{azulTitulo}\textbf{ELECTRÓNICA} --- Ejercicios}
\fancyhead[R]{\small\itshape Filtros Activos RC}
\fancyfoot[C]{\small Página \thepage\ de \pageref{LastPage}}
\renewcommand{\headrulewidth}{0.4pt}

% ── Colores y estilos específicos de ejercicios ────────────────────────────────
\definecolor{verdePuntaje}{HTML}{1A6B2F}

\tcbset{
  cajaInstrucciones/.style={
    enhanced, breakable,
    colback=grisFondo, colframe=grisBorde,
    fonttitle=\bfseries, coltitle=azulTitulo,
    top=6pt, bottom=6pt, left=8pt, right=8pt,
  },
}

% ──────────────────────────────────────────────────────────────────────────────
\begin{document}

% ══ PORTADA / ENCABEZADO INFográfico ════════════════════════════════════════════
\bandaTitulo{ELECTRÓNICA}{Filtros Activos RC}

\begin{center}
  \vspace{2pt}
  {\large \textbf{Dificultad:} Intermedia \hfill
   \textbf{Duración:} 90 minutos} \\[4pt]
  {\large \textbf{Fecha:} \rule{4cm}{0.2pt} \hfill
   \textbf{Estudiante:} \rule{6cm}{0.2pt}} \\[6pt]
  \rule{\linewidth}{0.2pt}
\end{center}

% ══ DATOS DE LOS EJERCICIOS ════════════════════════════════════════════════════
\vspace{4pt}
\begin{tcolorbox}[cajaInstrucciones, title={Instrucciones}]
Realice los cálculos de manera detallada y explique cada paso del proceso de resolución. Utilice fórmulas y ecuaciones de los filtros activos RC para obtener los resultados solicitados.
\medskip
\textbf{Material permitido:} Calculadora científica, formulario personal de una hoja.
\end{tcolorbox}

% ══ INDICACIÓN DE PUNTAJE ═════════════════════════════════════════════════════
\vspace{4pt}
\begin{center}
  \begin{tcolorbox}[
    enhanced,
    colback=white, colframe=azulTitulo,
    width=0.6\linewidth,
    boxrule=1.5pt,
    halign=center, valign=center,
    top=4pt, bottom=4pt,
  ]
    \centering
    \textbf{Puntaje total: 10.0 puntos} \\
    \small\textit{Responda cada ejercicio en el espacio provisto. \\
    Debe mostrar todo el desarrollo matemático para obtener puntaje completo.}
  \end{tcolorbox}
\end{center}

% ══ EJERCICIOS ═════════════════════════════════════════════════════════════════
\section*{Ejercicios}


% ── Ejercicio 1 ──────────────────────────────────────────────────────
\subsection*{Ejercicio 1 \normalfont\normalsize\textit{(2.0 puntos)}}
Diseñe un filtro pasa altos con una ganancia $A_v = 10$ y una frecuencia de corte $f_c = 1000 Hz$. La resistencia de entrada $R_i$ es $1 k\Omega$. ¿Cuáles son los valores de $R_1$, $R_2$ y $C$ necesarios para lograr las especificaciones del filtro?

  \begin{tcolorbox}[colback=gray!5, colframe=gray!40, fonttitle=\bfseries, title={Diagrama sugerido}]
    Un filtro pasa altos con amplificador operacional y componentes RC.
  \end{tcolorbox}

\vspace{2mm}
\rule{\linewidth}{0.2pt}
\vspace{3cm}  % espacio para respuesta



% ── Ejercicio 2 ──────────────────────────────────────────────────────
\subsection*{Ejercicio 2 \normalfont\normalsize\textit{(1.5 puntos)}}
Un filtro pasa bajas tiene una frecuencia de corte $f_c = 500 Hz$ y utiliza un condensador $C = 10 nF$. ¿Cuál es el valor de la resistencia $R$ necesaria para lograr esta frecuencia de corte?

  \begin{tcolorbox}[colback=gray!5, colframe=gray!40, fonttitle=\bfseries, title={Diagrama sugerido}]
    Un filtro pasa bajas con amplificador operacional y componentes RC.
  \end{tcolorbox}

\vspace{2mm}
\rule{\linewidth}{0.2pt}
\vspace{3cm}  % espacio para respuesta



% ── Ejercicio 3 ──────────────────────────────────────────────────────
\subsection*{Ejercicio 3 \normalfont\normalsize\textit{(2.5 puntos)}}
Diseñe un filtro pasa bandas con una frecuencia central $f_0 = 2000 Hz$, una banda pasante de $100 Hz$ y una ganancia $A_v = 5$. La resistencia de entrada $R_i$ es $2 k\Omega$. ¿Cuáles son los valores de los componentes necesarios para lograr las especificaciones del filtro?

  \begin{tcolorbox}[colback=gray!5, colframe=gray!40, fonttitle=\bfseries, title={Diagrama sugerido}]
    Un filtro pasa bandas con amplificador operacional y componentes RC.
  \end{tcolorbox}

\vspace{2mm}
\rule{\linewidth}{0.2pt}
\vspace{3cm}  % espacio para respuesta



% ── Ejercicio 4 ──────────────────────────────────────────────────────
\subsection*{Ejercicio 4 \normalfont\normalsize\textit{(2.0 puntos)}}
Un filtro notch tiene una frecuencia de resonancia $f_r = 1000 Hz$ y utiliza un condensador $C = 100 nF$. ¿Cuál es el valor de la resistencia $R$ necesaria para lograr esta frecuencia de resonancia?

  \begin{tcolorbox}[colback=gray!5, colframe=gray!40, fonttitle=\bfseries, title={Diagrama sugerido}]
    Un filtro notch con amplificador operacional y componentes RC.
  \end{tcolorbox}

\vspace{2mm}
\rule{\linewidth}{0.2pt}
\vspace{3cm}  % espacio para respuesta



% ── Ejercicio 5 ──────────────────────────────────────────────────────
\subsection*{Ejercicio 5 \normalfont\normalsize\textit{(2.5 puntos)}}
Diseñe un filtro rechaza bandas con una frecuencia central $f_0 = 1500 Hz$, una banda de rechazo de $200 Hz$ y una ganancia $A_v = 2$. La resistencia de entrada $R_i$ es $1 k\Omega$. ¿Cuáles son los valores de los componentes necesarios para lograr las especificaciones del filtro?

  \begin{tcolorbox}[colback=gray!5, colframe=gray!40, fonttitle=\bfseries, title={Diagrama sugerido}]
    Un filtro rechaza bandas con amplificador operacional y componentes RC.
  \end{tcolorbox}

\vspace{2mm}
\rule{\linewidth}{0.2pt}
\vspace{3cm}  % espacio para respuesta



% ══ FIN DE LOS EJERCICIOS ══════════════════════════════════════════════════════
\vspace{12pt}
\begin{center}
  \rule{0.6\linewidth}{0.4pt} \\[6pt]
  \textbf{--- Fin de la hoja de ejercicios ---}
\end{center}

\vfill
\begin{center}
  \footnotesize\color{gray}
  Ejercicios generados automáticamente por el sistema de inteligencia artificial (llama-3.3-70b-versatile) \\
  ELECTRÓNICA --- Sistema de Evaluación Académica \\
  Generado el 2026-06-27
\end{center}

\end{document}
